\section{Null moment fibers and positive duals which are not literal data}
\label{sec:r12-null-fibers-nonliteral-duals}

The first part uses precisely the proper original moment and
background factorization (ML12), with $1\leq q<r$:
\[
 t:V\longrightarrow Q\subset\mathbf R^r,\qquad
 S=(t^*\beta_0)^{r-q}\wedge R,\qquad
 \beta_0=\sum_{i=1}^r d'x_i\wedge d''x_i.
 \tag{NF1}
\]
Here $R$ is an arbitrary allowed original symmetric middle-positive
separately closed $(n-r,n-r)$-current. Properness means that the
inverse image of every compact real set is compact in the ordinary
analytic open $V$. All capacities use every normalized original
bounded classical tester on $V$, without restricting its actual
moment refinements. The factorization is a printed hypothesis for
these geometric null-set conclusions, not a property of every $S$.

\begin{lemma}[A closed order-zero supercurrent has no sub-top atom]
\label{lem:r12-closed-current-no-atom}
Let $T$ be an order-zero real supercurrent of bidegree $(q,q)$
on $Q$, with $q<r$ and $d'T=0$. Every coefficient measure of
$T$ has zero mass at each point. In particular, if $T$ is
middle-positive, the positive measure
$T\wedge\beta_0^{r-q}$ has no atoms.
\end{lemma}

\noindent\textbf{Status: \texttt{DIRECT\_COROLLARY}} of
distributional closedness and the coefficient-measure definition;
the elementary proof is included. Separate $d''$-closedness and
positivity are not needed for the first assertion.

\begin{proof}
Fix $x\in Q$ and write the vector of coefficient measures near
$x$ as its atom at $x$ plus a vector of measures having no atom
there. Let $k=r-q\geq1$, let $\eta$ be any constant form of
degree $(k-1,k)$, and choose $\varphi\in C_c^\infty(\mathbf R^r)$.
For small $\epsilon>0$ the compact test
\[
               \varphi_\epsilon(y)\eta
           =\epsilon\varphi((y-x)/\epsilon)\eta
\]
is supported in $Q$. Closedness gives
$\langle T,d'\varphi_\epsilon\wedge\eta\rangle=0$.
The coefficients of $d'\varphi_\epsilon$ are uniformly bounded.
Thus the contribution from the atom-free part tends to zero,
because its variations on these shrinking balls tend to zero.
The contribution of the atom is its pairing with
$d'\varphi(0)\wedge\eta$. Every covector can occur as
$d'\varphi(0)$, and forms $\xi'\wedge\eta$ of these degrees
span the entire complementary coefficient space $(k,k)$.
The atomic tensor therefore pairs to zero with every coefficient
test and is zero. This works at every point.
The trace is a fixed finite linear combination of coefficient
measures, so it has no atom either.
\end{proof}

\begin{theorem}[Full original outer capacity vanishes on each proper moment fiber]
\label{thm:r12-null-proper-moment-fibers}
Under (NF1), let $H\Subset Q$ be compact. For sufficiently
small $\epsilon>0$ put
\[
 \omega_H(\epsilon)=
   \sup_{x\in H}\ \sup_{\substack{v\text{ original bounded classical psh on }V\\
                                  0\leq v\leq1}}
     \bigl((Dv)^q\wedge S\bigr)
                           \bigl(t^{-1}(B(x,\epsilon))\bigr).
 \tag{NF2}
\]
This is finite and
\[
                  \lim_{\epsilon\downarrow0}\omega_H(\epsilon)=0.
 \tag{NF3}
\]
Consequently, for every $x\in Q$,
\[
                    C_S^*(t^{-1}\{x\};V)=0.
 \tag{NF4}
\]
The same holds for the inverse image of any countable subset
of $Q$. The conclusions concern the full original capacity
and ordinary open neighborhoods, not a toric-only tester capacity.
No numerical power rate is asserted in (NF3).
\end{theorem}

\noindent\textbf{Status: \texttt{NEW\_PROVED}} as a
uniform full-tester null-fiber statement in every sub-top
degree. It uses the new positive moment lifts and the frozen
arbitrary-background relative CLN estimate.

\begin{proof}
For every tester in (NF2), its original bounded lift is
\[
             T_v=t_*\bigl((Dv)^q\wedge R\bigr),\qquad
             \nu_v=T_v\wedge\beta_0^{r-q}
                       =t_*\bigl((Dv)^q\wedge S\bigr).
 \tag{NF5}
\]
These are the actual currents and exact unnormalized trace
identity of (ML15)--(ML17). On each compact $H'\Subset Q$,
properness makes $t^{-1}(H')$ compact. The original relative
CLN estimate on a slightly larger compact set, with
$\|v\|_\infty\leq1$, bounds $\nu_v(H')$ uniformly over
the full tester class. This proves finiteness in (NF2).
The trace coefficient lemma (ML4)--(ML5) makes the family
of lifts locally sequentially compact. Every limiting lift
is again positive, symmetric and separately closed.

If (NF3) failed, there would be $\delta>0$,
$\epsilon_j\downarrow0$, $x_j\in H$, and normalized
testers $v_j$ with
$\nu_{v_j}(B(x_j,\epsilon_j))\geq\delta$.
After subsequences, $x_j\to x\in H$ and $T_{v_j}\to T$
locally on all real tests. Its trace $\nu$ is the local vague
limit of $\nu_{v_j}$. For any sufficiently small fixed
$\rho>0$, eventually
$B(x_j,\epsilon_j)\subset\overline B(x,\rho)\Subset Q$.
Portmanteau on this compact set gives
\[
      \delta\leq\limsup_j\nu_{v_j}(\overline B(x,\rho))
                              \leq\nu(\overline B(x,\rho)).
\]
Letting $\rho\downarrow0$ gives $\nu\{x\}\geq\delta$,
contradicting Lemma~\ref{lem:r12-closed-current-no-atom}.
This proves the uniform assertion.

The sets $t^{-1}(B(x,\epsilon))$ are ordinary open
neighborhoods of the compact fiber $t^{-1}\{x\}$.
Their raw capacities are exactly the inner supremum in
(NF2). Taking their decreasing numerical bound in the
definition of outer capacity proves (NF4). Countable
subadditivity of the frozen outer capacity gives the last
assertion. Neither a Radon disintegration over residue
fields nor a change of the original field has occurred.
\end{proof}

\subsection{Canonical convex traces of the complete energy class}

For the remainder take the allowed original example
\[
 \begin{gathered}
 X=(\mathbf G_{m,K}^{r})^{\rm an},\quad
 Q=(-1,1)^r,\quad V=t^{-1}(Q),\quad t_i=\log|T_i|,\\
 1\leq q<r,\quad S=\beta^{r-q}\wedge[V],\quad
 \beta=\sum_i d't_i\wedge d''t_i,\quad
 \sigma:Q\longrightarrow V\text{ the original joint Gauss section}.
 \end{gathered}
 \tag{NF6}
\]
These data satisfy (ED1) and (NF1). The frozen R6
\texttt{lem:r6-joint-gauss-convexity}, used in R11
\texttt{comparison/power\_class\_source\_obstruction.tex}
before the choice of its background, says that a bounded
original classical psh potential has a finite convex
restriction to $\sigma(Q)$. That assertion is independent
of the remaining degree $q$. We use it as an exact
\texttt{INTERNAL\_INPUT}; no nonlinear identity between
a Gauss trace Hessian and the original top measure is used.

\begin{proposition}[All-degree completed Gauss traces and point evaluation]
\label{prop:r12-completed-gauss-evaluation}
Every $u\in\mathscr C$ for (NF6) has a canonical finite
nonpositive continuous convex function $g_u$ on $Q$.
It is the locally uniform limit of the Gauss traces of any
actual bounded strong defining row. It agrees with the
chosen relative Borel representative $u\circ\sigma$ only
Lebesgue-almost everywhere, not necessarily at every point.

The assignment is additive and positively homogeneous on
$\mathscr C$, and extends linearly to $\mathbb D_S$.
For each $x\in Q$, the formula
\[
                         \ell_x(u)=-g_u(x)
 \tag{NF7}
\]
defines an order-reversing functional in
$\mathbb D_{S,\mathrm s}^*$. For every original compact
smooth scalar $\phi$ on $V$,
\[
                              \ell_x(\phi)=-\phi(\sigma(x)).
 \tag{NF8}
\]
Nevertheless this functional is not literal integration
against a Radon measure on all relative Borel representatives
in $\mathbb D_S$.
\end{proposition}

\noindent\textbf{Status: \texttt{NEW\_PROVED}} for the
full-degree strong class and the nonliteral positive dual
phenomenon. The old bounded Gauss convexity is not recounted
as a new result.

\begin{proof}
Fix real compact sets $H\Subset\operatorname{Int}H'\Subset Q$.
Patch the quadratic $|t|^2/2-C$ to a large multiple of the
strict-box barrier, as in (ML25), to obtain an actual
bounded first-energy potential $\psi$ equal to that quadratic
on $t^{-1}(H')$. All factors are unnormalized, so locality
and the original torus calibration give
\[
           \mathbf1_{t^{-1}(H')}\mu_\psi
                   =\mathbf1_{t^{-1}(H')}\beta^r\wedge[V]
                   =r!\,\sigma_*(\mathbf1_{H'}\,dx).
 \tag{NF9}
\]
For any actual bounded first-energy potential $a$,
the mixed energy inequality implies
\[
              \int_{H'}(-a\circ\sigma)\,dx
                          \leq\frac{r(\psi)^q}{r!}\,r(a).
 \tag{NF10}
\]
For two such potentials, the completed tester estimate
(EV1) controls their trace difference in $L^1(H')$ by
a quantity tending to zero under strong convergence on
an energy-bounded set.

We recall explicitly the elementary convex compactness step.
If $g\leq0$ is finite convex on a ball $B(x,2\rho)$,
then $-g$ is nonnegative and concave. For
$y=(x+z)/2\in B(x,\rho)$, concavity gives
$-g(y)\geq(-g(x))/2$. Hence
$-g(x)\leq2\operatorname{vol}(B(x,\rho))^{-1}
\int_{B(x,\rho)}(-g)$. An $L^1$ bound on a larger
compact set therefore bounds $|g|$ on an interior compact
set. Secant-slope estimates for convex functions then give
a uniform Lipschitz bound on a still smaller compact set.
Thus an $L^1_{\rm loc}$-Cauchy row of these convex functions
converges locally uniformly to one finite continuous convex
function: compactness gives subsequential limits, and its
$L^1$ limit makes them identical.

Choose an actual bounded strongly defining row for $u$;
such a row is supplied by the R11 clipping diagonal.
The traces are convex by the frozen input, are locally
$L^1$-bounded by (NF10), and are $L^1$-Cauchy by (EV1).
The preceding argument defines $g_u$ and proves independence
of the row. Their $L^1(\mu_\psi)$ convergence to the
relative Borel class, using (NF9), proves the almost-everywhere
identification in the statement. Exact cone addition and
homogeneity, followed by local $L^1$ limits, give the stated
linearity. Equal Borel difference presentations give equal
traces almost everywhere, hence everywhere by continuity.

The point bound from (NF10) yields, for each fixed $x$,
\[
                     0\leq-g_u(x)\leq C_x r(u),
 \tag{NF11}
\]
with a finite constant independent of $u$. Extending by
differences and infimizing proves boundedness of $\ell_x$
on $\mathbb D_S$. If $h=a-b\succeq0$, its relative
exceptional set is $\mu_\psi$-null. Thus
$g_a-g_b\geq0$ almost everywhere on each core, and continuity
makes the inequality hold at every point. This proves order
reversal. For a strongly convergent cone row, (EV1) and
(NF9) again give local $L^1$ convergence of its convex
traces; the uniform convex bounds give local uniform
convergence. Consequently $\ell_x$ is strongly continuous.

The all-degree original test decomposition (HC14)--(HC15)
writes $\phi=a-b$ as an exact difference of actual bounded
finite-mass common-collar potentials. Their actual Gauss
traces give (NF8), including every refined non-toric original
test. In particular any Radon measure giving the same
original test restriction must be $\delta_{\sigma(x)}$.
But (NF4) makes $\{\sigma(x)\}$ outer-capacity null.
Changing a Borel representative at this point does not
change its class, whereas literal integration against this
Dirac measure changes. For example, change the value of
the strictly negative patched $\psi$ there to zero.
Thus the functional cannot be literal Radon integration
on all representatives of $\mathbb D_S$. This does not
conflict with its canonical-trace definition.
\end{proof}

\begin{theorem}[Literal positive data are not closed in the energy dual norm]
\label{thm:r12-literal-dual-data-not-closed}
In (NF6), for every $x\in Q$ there are finite positive Radon
measures $\mu_\epsilon$, each dominated by a finite constant
times one actual bounded first-energy source, such that
their literal integration functionals satisfy
\[
               \|\ell_{\mu_\epsilon}-\ell_x\|_*
                                      \leq C_x'\epsilon\longrightarrow0.
 \tag{NF12}
\]
All these functionals and their limit are order reversing
and strongly continuous. The limit is not a literal Radon
datum on the relative Borel quotient. The same measures
converge on every original compact smooth scalar test to
$\delta_{\sigma(x)}$, which charges a full outer-capacity
null set.

This does not contradict closedness of $J(\mathscr C)$:
the approximating data here are not source functionals.
Their nonzero compact tropical carriers already exclude
that possibility by (ML20).
\end{theorem}

\noindent\textbf{Status: \texttt{NEW\_PROVED}} as a
counterexample to the blanket identification of the positive
strong energy dual with literal finite-energy Radon data.
That identification, and norm-closedness of the literal
positive datum class without a uniform diffuseness gate,
are \texttt{FALSE}.

\begin{proof}
Choose a nonnegative smooth probability density
$\theta_\epsilon$ supported in $B(x,\epsilon)\Subset Q$
and set $\mu_\epsilon=\sigma_*(\theta_\epsilon\,dx)$.
On one fixed larger compact core, (NF9) gives
$\mu_\epsilon\leq(\|\theta_\epsilon\|_\infty/r!)\mu_\psi$.
The mixed energy inequality and (EV1) therefore make
each a literal member of $\mathbb D_{S,\mathrm s}^*$.
For $u\in\mathscr C$ with $r(u)\leq1$, the proof of
(NF11) also gives a common Lipschitz constant for $g_u$
on this core. The almost-everywhere identification of the
representative gives
\[
 \left|\ell_{\mu_\epsilon}(u)-\ell_x(u)\right|
       =\left|\int(g_u(x)-g_u(y))\theta_\epsilon(y)\,dy\right|
       \leq C_x'\epsilon.
\]
The exact cone formula for the operator dual norm (ED10)
proves (NF12). Evaluation on continuous original compact
tests proves the ordinary measure convergence. The previous
proposition proves that its limit is nonliteral on the
relative quotient. Finally the carrier of each tropical
pushforward is a nonempty compact subset of $Q$, so the
source obstruction applies to every approximating measure.
All constructions remain on the original field and norm.
\end{proof}

\paragraph{The exact surviving distinctions.}
There are now three separate notions: a bounded Banach-dual
functional, a strong-continuous order-reversing functional,
and a literal positive Radon datum whose integrals ignore
the relative null ideal. The first need not be strongly
continuous by (CG11)--(CG14), and the second need not be
literal by (NF7)--(NF12). The capacity-modulus condition
(CD1)--(CD2) supplies the third with all required continuity.
The nonlinear source image is a further, closed proper subset,
with the normal and contact obstructions already proved.
None of these distinctions closes OPEN-GQ or identifies
relative null sets with all pluripolar or negligible sets.
